\section*{Adaptaciones a CRISP-DM}
\addcontentsline{toc}{section}{Adaptaciones a CRISP-DM}

En esta fase hicimos los siguientes ajustes a CRISP-DM:

\begin{itemize}
  \item Mantenemos el orden de las tareas 3.1 a 3.5, pero la conversión de las fotos a JPG, que es parte de 3.5, la hacemos antes de 3.3. La vista de cada foto se marca a mano en una página que corre en el navegador, y el navegador no puede mostrar fotos HEIC.
  \item Preparamos solo los datos de BCS (fotos y bitácora), porque los datos de ordeño no son entrada del modelo. Por eso esta fase usa un solo notebook de Python. Se llama \nolinkurl{01_dp_bcs_problem.ipynb} y está en la carpeta \nolinkurl{src/notebooks/03-data-preparation/} del repositorio.
  \item Todo se escribe en la carpeta \nolinkurl{data/}, que no se sube al repositorio, y los archivos originales nunca se modifican.
  \item La separación en 5 pliegues por vaca es la tarea 4.2 de CRISP-DM (diseño de prueba), pero la hacemos aquí para que el dataset salga listo para entrenar.
  \item No hay registros generados. La descripción del dataset la hacemos con las tablas de columnas de la sección~\ref{sec:descripcion} y con un \nolinkurl{LEEME.md} que mantiene el equipo en el bucket de S3.
\end{itemize}
